\section{Análisis de Complejidad Teórica}

\subsection{Representación: Lista de Adyacencia}

Se utiliza una lista de adyacencia implementada con \texttt{unordered\_map<int, vector<Edge>>}.
Esta representación es óptima para redes reales, que son típicamente dispersas ($|E| \ll |V|^2$)~\cite{cormen2009}.

\begin{center}
\begin{tabular}{lccc}
\toprule
\textbf{Representación} & \textbf{Memoria} & \textbf{Iterar vecinos} & \textbf{Verificar arista} \\
\midrule
Lista de adyacencia & $O(|V|+|E|)$ & $O(\deg(v))$ & $O(\deg(v))$ \\
Matriz de adyacencia & $O(|V|^2)$ & $O(|V|)$ & $O(1)$ \\
\bottomrule
\end{tabular}
\end{center}

Para Yeast con $|V| = 6.526$, una matriz requeriría $6.526^2 \approx 42$ millones de
entradas ($\sim 340$ MB), mientras que la lista de adyacencia ocupa $\sim 24$ MB.

\subsection{Operaciones sobre el Grafo}

\begin{center}
\begin{tabular}{lc}
\toprule
\textbf{Operación} & \textbf{Complejidad} \\
\midrule
\texttt{addVertex} & $O(1)$ esperado \\
\texttt{addEdge} & $O(1)$ esperado \\
\texttt{removeEdge} & $O(\deg(v))$ \\
\texttt{hasEdge} & $O(\deg(v))$ \\
Construcción completa & $O(|V|+|E|)$ \\
\bottomrule
\end{tabular}
\end{center}

\texttt{addVertex} inserta en \texttt{unordered\_map} y \texttt{addEdge} hace una búsqueda
en \texttt{unordered\_map} más un \texttt{push\_back} en \texttt{vector}. Ambas son $O(1)$
\textbf{esperado} (caso promedio con función de hash uniforme), no amortizado en el sentido
estricto: en el peor caso con colisiones degeneradas, la complejidad sube a $O(|V|)$.

\subsection{Complejidad de las Métricas}

\begin{center}
\begin{tabular}{lcc}
\toprule
\textbf{Métrica} & \textbf{Complejidad} & \textbf{Algoritmo base} \\
\midrule
Degree Centrality & $O(|V|+|E|)$ & Recorrido de vecinos \\
Betweenness Centrality & $O(|V| \cdot |E|)$ & Brandes (BFS por fuente) \\
Closeness Centrality & $O(|V| \cdot (|V|+|E|))$ & BFS por fuente \\
PageRank & $O(k \cdot (|V|+|E|))$ & Iteración de potencias \\
Average Shortest Path & $O(|V| \cdot (|V|+|E|))$ & BFS por fuente \\
Eigenvector Centrality & $O(k \cdot (|V|+|E|))$ & Power Iteration \\
Clustering Coefficient & $O(|V| \cdot \overline{\deg}^{\,2})$ & Conteo de triángulos \\
\bottomrule
\end{tabular}
\end{center}

donde $k$ es el número de iteraciones hasta convergencia ($k < 100$ en la práctica)
y $\overline{\deg}$ es el grado promedio de la red.

\textbf{Peor caso dominante:} Betweenness con $O(|V| \cdot |E|)$. Para Yeast esto
implica $6.526 \times 532.180 \approx 3{,}47 \times 10^9$ operaciones por ejecución,
lo que se verá reflejado en los experimentos.

\subsection{Relación entre Métricas}

Las métricas no son independientes. En redes reales existe correlación positiva entre
Degree y Eigenvector Centrality (hubs se conectan con hubs), mientras que Betweenness
puede identificar vértices con degree moderado pero posición estructural estratégica
(puentes entre comunidades). Esta comparación es parte del análisis experimental.
